\textbf{结论名称：}常见不等式的证明方法汇总（比较法/综合法/分析法/反证法）

\textbf{适用条件：}比较法适用于可直接作差或作商变形的情形；综合法适用于已知条件明确、可逐步推导出结论的情形；分析法适用于从结论出发逆向寻找充分条件；反证法适用于直接证明困难、结论为否定形式或唯一性命题.

\textbf{变量说明：} $a,b$ 为实数或代数式；$P, Q$ 表示一般的命题.

\textbf{核心公式：}
比较法的基本等价形式：
\[
\boxed{a>b \iff a-b>0,\qquad a=b \iff a-b=0,\qquad a<b \iff a-b<0}
\]
综合法与分析法互为逆过程：
\[
\text{综合法（由因导果）：}\; P \Rightarrow Q_1 \Rightarrow Q_2 \Rightarrow \cdots \Rightarrow Q;
\]
\[
\text{分析法（执果索因）：}\; Q \Leftarrow P_1 \Leftarrow P_2 \Leftarrow \cdots \Leftarrow P.
\]
反证法的逻辑结构：
\[
\neg Q \land P \Rightarrow \text{矛盾} \;\Longrightarrow\; P \Rightarrow Q.
\]

\textbf{使用场景：}证明不等式恒成立、比较代数式大小、证明唯一性或否定性命题.

\textbf{简单例子：}证明 $a^2+b^2 \ge 2ab$（$a,b\in\mathbb{R}$）：
比较法——作差：$a^2+b^2-2ab = (a-b)^2 \ge 0$，故 $a^2+b^2 \ge 2ab$；
综合法——由 $(a-b)^2 \ge 0$ 展开得 $a^2-2ab+b^2 \ge 0$，移项即得 $a^2+b^2 \ge 2ab$；
分析法——要证 $a^2+b^2 \ge 2ab$，只需证 $a^2-2ab+b^2 \ge 0$，即 $(a-b)^2 \ge 0$，此式显然成立；
反证法——假设存在 $a,b$ 使 $a^2+b^2 < 2ab$，则 $(a-b)^2 < 0$，与平方非负矛盾，故原不等式成立.